\documentclass[aspectratio=169]{beamer}

\usetheme{Madrid}
\usecolortheme{default}

\usepackage[spanish,es-noshorthands]{babel}
\usepackage[utf8]{inputenc}
\usepackage[T1]{fontenc}
\usepackage{lmodern}
\usepackage{amsmath,amssymb}
\usepackage{array}
\usepackage{xcolor}

\definecolor{USTGreen}{RGB}{27,94,32}
\definecolor{USTLight}{RGB}{232,245,233}

\setbeamercolor{structure}{fg=USTGreen}
\setbeamercolor{frametitle}{fg=white,bg=USTGreen}
\setbeamercolor{title}{fg=white,bg=USTGreen}
\setbeamertemplate{navigation symbols}{}

\title[Sistemas de ecuaciones lineales]{Sistemas de Ecuaciones Lineales con Dos Incógnitas}
\subtitle{Clasificación, métodos algebraicos e interpretación}
\author{Fundamentos de Matemática Básica}
\institute{Universidad Santo Tomás}
\date{}

\begin{document}
	
	%------------------------------------------------
	\begin{frame}
		\titlepage
	\end{frame}
	
	%------------------------------------------------
	\begin{frame}{Objetivo general}
		Analizar y resolver sistemas de ecuaciones lineales con dos incógnitas mediante métodos algebraicos y geométricos, interpretando el significado de sus soluciones.
		
		\bigskip
		
		\begin{block}{Resultados de aprendizaje}
			\begin{itemize}
				\item Identificar sistemas compatibles determinados, compatibles indeterminados e incompatibles.
				\item Resolver sistemas mediante sustitución, igualación y reducción.
				\item Interpretar geométricamente la solución de un sistema.
				\item Detectar patrones, relaciones e inconsistencias.
			\end{itemize}
		\end{block}
	\end{frame}
	
	%------------------------------------------------
	\begin{frame}{Definición}
		Un sistema de ecuaciones lineales con dos incógnitas tiene la forma:
		
		\[
		\begin{cases}
			ax+by=c\\
			dx+ey=f
		\end{cases}
		\]
		
		donde:
		
		\[
		a,b,c,d,e,f\in\mathbb{R}
		\]
		
		\bigskip
		
		Resolver el sistema significa encontrar los valores de \(x\) e \(y\) que satisfacen simultáneamente ambas ecuaciones.
	\end{frame}
	
	%------------------------------------------------
	\begin{frame}{Interpretación geométrica}
		Cada ecuación lineal representa una recta en el plano cartesiano.
		
		\bigskip
		
		Por lo tanto, resolver un sistema equivale a estudiar la relación entre dos rectas.
		
		\bigskip
		
		\begin{center}
			\begin{tabular}{|c|c|c|}
				\hline
				\textbf{Caso} & \textbf{Rectas} & \textbf{Soluciones} \\
				\hline
				Compatible determinado & Se cortan & Una solución \\
				\hline
				Compatible indeterminado & Coinciden & Infinitas soluciones \\
				\hline
				Incompatible & Paralelas & Sin solución \\
				\hline
			\end{tabular}
		\end{center}
	\end{frame}
	
	%------------------------------------------------
	\begin{frame}{Sistema compatible determinado}
		Tiene una única solución.
		
		\[
		\begin{cases}
			x+y=5\\
			2x-y=4
		\end{cases}
		\]
		
		\bigskip
		
		Sumamos ambas ecuaciones:
		
		\[
		(x+y)+(2x-y)=5+4
		\]
		
		\[
		3x=9
		\]
		
		\[
		x=3
		\]
		
		Luego:
		
		\[
		3+y=5
		\]
		
		\[
		y=2
		\]
		
		\[
		\boxed{(x,y)=(3,2)}
		\]
	\end{frame}
	
	%------------------------------------------------
	\begin{frame}{Sistema compatible indeterminado}
		Tiene infinitas soluciones.
		
		\[
		\begin{cases}
			x+y=4\\
			2x+2y=8
		\end{cases}
		\]
		
		La segunda ecuación es múltiplo de la primera:
		
		\[
		2(x+y)=8
		\]
		
		\[
		x+y=4
		\]
		
		\bigskip
		
		Ambas ecuaciones representan la misma recta.
		
		\[
		\boxed{\text{Infinitas soluciones}}
		\]
	\end{frame}
	
	%------------------------------------------------
	\begin{frame}{Sistema incompatible}
		No tiene solución.
		
		\[
		\begin{cases}
			x+y=3\\
			x+y=7
		\end{cases}
		\]
		
		\bigskip
		
		La misma expresión \(x+y\) no puede ser igual a 3 y a 7 al mismo tiempo.
		
		\bigskip
		
		Esto genera una contradicción:
		
		\[
		3=7
		\]
		
		\[
		\boxed{\text{Sistema incompatible}}
		\]
	\end{frame}
	
	%------------------------------------------------
	\begin{frame}{Método de sustitución}
		Consiste en despejar una incógnita en una ecuación y reemplazarla en la otra.
		
		\bigskip
		
		Resolver:
		
		\[
		\begin{cases}
			x+y=10\\
			x=y+2
		\end{cases}
		\]
		
		Sustituimos \(x=y+2\):
		
		\[
		(y+2)+y=10
		\]
		
		\[
		2y+2=10
		\]
		
		\[
		2y=8
		\]
		
		\[
		y=4
		\]
		
		\[
		x=4+2=6
		\]
		
		\[
		\boxed{(x,y)=(6,4)}
		\]
	\end{frame}
	
	%------------------------------------------------
	\begin{frame}{Método de igualación}
		Consiste en despejar la misma incógnita en ambas ecuaciones e igualar las expresiones.
		
		\bigskip
		
		Resolver:
		
		\[
		\begin{cases}
			x=8-y\\
			x=2y+2
		\end{cases}
		\]
		
		Igualamos:
		
		\[
		8-y=2y+2
		\]
		
		\[
		6=3y
		\]
		
		\[
		y=2
		\]
		
		Luego:
		
		\[
		x=8-2=6
		\]
		
		\[
		\boxed{(x,y)=(6,2)}
		\]
	\end{frame}
	
	%------------------------------------------------
	\begin{frame}{Método de reducción}
		Consiste en eliminar una incógnita sumando o restando ecuaciones.
		
		\bigskip
		
		Resolver:
		
		\[
		\begin{cases}
			2x+y=11\\
			x-y=1
		\end{cases}
		\]
		
		Sumamos:
		
		\[
		(2x+y)+(x-y)=11+1
		\]
		
		\[
		3x=12
		\]
		
		\[
		x=4
		\]
		
		Sustituimos:
		
		\[
		4-y=1
		\]
		
		\[
		y=3
		\]
		
		\[
		\boxed{(x,y)=(4,3)}
		\]
	\end{frame}
	
	%------------------------------------------------
	\begin{frame}{Ejercicio 1}
		Resuelva por sustitución:
		
		\[
		\begin{cases}
			x+y=18\\
			x=2y
		\end{cases}
		\]
		
		\bigskip
		
		\begin{block}{Preguntas}
			\begin{itemize}
				\item ¿Cuál es la solución del sistema?
				\item ¿Qué tipo de sistema es?
				\item ¿Cómo se verifica la solución?
			\end{itemize}
		\end{block}
	\end{frame}
	
	%------------------------------------------------
	\begin{frame}{Solución ejercicio 1}
		Sistema:
		
		\[
		\begin{cases}
			x+y=18\\
			x=2y
		\end{cases}
		\]
		
		Sustituimos:
		
		\[
		2y+y=18
		\]
		
		\[
		3y=18
		\]
		
		\[
		y=6
		\]
		
		\[
		x=2(6)=12
		\]
		
		\[
		\boxed{(x,y)=(12,6)}
		\]
		
		Es un sistema compatible determinado.
	\end{frame}
	
	%------------------------------------------------
	\begin{frame}{Ejercicio 2}
		Resuelva por igualación:
		
		\[
		\begin{cases}
			x=20-2y\\
			x=y+5
		\end{cases}
		\]
		
		\bigskip
		
		\begin{block}{Tarea}
			Resolver, clasificar e interpretar geométricamente.
		\end{block}
	\end{frame}
	
	%------------------------------------------------
	\begin{frame}{Solución ejercicio 2}
		Igualamos:
		
		\[
		20-2y=y+5
		\]
		
		\[
		15=3y
		\]
		
		\[
		y=5
		\]
		
		Luego:
		
		\[
		x=5+5=10
		\]
		
		\[
		\boxed{(x,y)=(10,5)}
		\]
		
		El sistema es compatible determinado.
	\end{frame}
	
	%------------------------------------------------
	\begin{frame}{Ejercicio 3}
		Resuelva por reducción:
		
		\[
		\begin{cases}
			3x+2y=24\\
			x-2y=0
		\end{cases}
		\]
		
		\bigskip
		
		\begin{block}{Tarea}
			Resolver el sistema y verificar la solución.
		\end{block}
	\end{frame}
	
	%------------------------------------------------
	\begin{frame}{Solución ejercicio 3}
		Sumamos ambas ecuaciones:
		
		\[
		(3x+2y)+(x-2y)=24+0
		\]
		
		\[
		4x=24
		\]
		
		\[
		x=6
		\]
		
		Sustituimos:
		
		\[
		6-2y=0
		\]
		
		\[
		2y=6
		\]
		
		\[
		y=3
		\]
		
		\[
		\boxed{(x,y)=(6,3)}
		\]
	\end{frame}
	
	%------------------------------------------------
	\begin{frame}{Patrones e inconsistencias}
		Durante la resolución pueden aparecer patrones importantes:
		
		\bigskip
		
		\begin{block}{Contradicción}
			Si aparece una igualdad falsa:
			
			\[
			0=5
			\]
			
			el sistema es incompatible.
		\end{block}
		
		\begin{block}{Identidad}
			Si aparece una igualdad verdadera:
			
			\[
			0=0
			\]
			
			el sistema es compatible indeterminado.
		\end{block}
	\end{frame}
	
	%------------------------------------------------
	\begin{frame}{Actividad de clasificación}
		Clasifique cada sistema:
		
		\[
		A)
		\begin{cases}
			x+y=10\\
			2x+2y=20
		\end{cases}
		\]
		
		\[
		B)
		\begin{cases}
			x+y=10\\
			x+y=15
		\end{cases}
		\]
		
		\[
		C)
		\begin{cases}
			x+y=10\\
			x-y=2
		\end{cases}
		\]
	\end{frame}
	
	%------------------------------------------------
	\begin{frame}{Solución actividad de clasificación}
		\[
		A)
		\begin{cases}
			x+y=10\\
			2x+2y=20
		\end{cases}
		\]
		
		Compatible indeterminado.
		
		\bigskip
		
		\[
		B)
		\begin{cases}
			x+y=10\\
			x+y=15
		\end{cases}
		\]
		
		Incompatible.
		
		\bigskip
		
		\[
		C)
		\begin{cases}
			x+y=10\\
			x-y=2
		\end{cases}
		\]
		
		Compatible determinado.
	\end{frame}
	
	%------------------------------------------------
	\begin{frame}{Cierre}
		\begin{block}{Ideas clave}
			\begin{itemize}
				\item Un sistema lineal con dos incógnitas representa dos rectas.
				\item Puede tener una solución, infinitas soluciones o ninguna solución.
				\item Los métodos algebraicos principales son sustitución, igualación y reducción.
				\item La solución debe verificarse algebraicamente y analizarse en contexto.
			\end{itemize}
		\end{block}
	\end{frame}
	
	%------------------------------------------------
	\begin{frame}{Tarea}
		Resolver y clasificar los siguientes sistemas:
		
		\[
		1)
		\begin{cases}
			2x+y=15\\
			x-y=3
		\end{cases}
		\]
		
		\[
		2)
		\begin{cases}
			x+y=8\\
			3x+3y=24
		\end{cases}
		\]
		
		\[
		3)
		\begin{cases}
			2x-y=5\\
			4x-2y=12
		\end{cases}
		\]
		
		\bigskip
		
		Indicar si el sistema es compatible determinado, compatible indeterminado o incompatible.
	\end{frame}
	
\end{document}